% !TeX program = xelatex
\documentclass[a4paper,runningheads]{llncs}
\usepackage[UTF8,scheme=plain,fontset=fandol]{ctex}
\usepackage{amsmath,booktabs,tabularx,graphicx,xcolor,listings,url}
\usepackage[hidelinks]{hyperref}
\lstdefinestyle{code}{basicstyle=\ttfamily\footnotesize,columns=fullflexible,frame=single,numbers=left,numberstyle=\scriptsize\color{gray},breaklines=true,showstringspaces=false,keepspaces=true,xleftmargin=1.4em,framexleftmargin=1.1em}
\lstset{style=code}
\newcommand{\code}[1]{\texttt{#1}}
\title{深入理解编译系统：一个 SysY 程序从源代码到 RISC-V 可执行文件}
\titlerunning{一个 SysY 程序的编译系统纵向解剖}
\author{成员甲（姓名与学号）\inst{1} \and 成员乙（姓名与学号）\inst{1}}
\authorrunning{成员甲与成员乙}
\institute{编译系统原理课程实验组}
\emergencystretch=2em
\begin{document}
\maketitle
\begin{abstract}
“编译器把源代码变成机器码”掩盖了各工具之间最重要的接口。本文把一个有界迭代阶乘 SysY 程序作为同一条语义线索，依次考察预处理文本、带类型抽象语法树、LLVM 中间表示、RV64GC 汇编、ELF 可重定位对象、静态可执行文件与 Linux 进程映像。核心观点是：编译系统是一组\emph{表示契约}；每个边界既遗忘已经无用的源形式，又显式加入下一消费者所需的事实。实例把赋值和循环追踪为内存读写与静态单赋值数据依赖，把函数调用追踪为 RV64 psABI 寄存器约定、符号和重定位，再把链接后的节视图追踪为加载时的段视图。Clang 18.1.3 的观测显示，\code{-O0} 与 \code{-O2} 产物在结构上显著不同，但结构计数不等于性能。源程序、作者书写的 LLVM IR 和 RV64 汇编在六个边界输入上均为 6/6 等价；这是针对声明输入的行为检查而非形式证明。最后以 MLIR 的多级方言和渐进式降低说明：编译设计的关键不是尽快降低，而是在最后一个需要某种抽象的变换完成后再丢弃它。
\keywords{编译系统 \and 表示契约 \and LLVM IR \and 静态单赋值 \and RISC-V \and ELF \and MLIR}
\end{abstract}
\section*{作者分工}
成员甲负责前端与 LLVM IR 的材料整理和初稿；成员乙负责 RV64、ELF、MLIR 展望与排版。双方共同确定案例语义、复核六个输入的行为、整理文献并终审全文。
\section{引言：不要把编译器画成一个黑箱}
\label{sec:introduction}

在命令行输入一次 \texttt{clang program.c}，得到一个可以运行的文件，很容易造成一种过于
粗糙的心智模型：源程序进入一个名为“编译器”的黑箱，机器程序从另一端出来。这个模型
没有错到不能使用，却粗糙到无法回答任何真正的系统问题。例如，宏为何不会成为抽象语法
树中的节点？局部变量何时不再对应一块内存？外部函数的地址在目标文件中尚未确定，调用
指令又如何能够先被编码？优化后的程序表面上已面目全非，我们依据什么称其“仍是同一个
程序”？这些问题分别属于不同的表示、不同的责任边界，也需要不同种类的证据。

本文采用一个更有解释力的观点：编译系统是一条\emph{表示契约}
（representation contract）链。每个阶段接收一种表示，确立或利用其中的事实，把其中一部分
决策显式化，同时把尚不能决定的问题连同足够的信息交给下一阶段。这里的“契约”不是指每个
内部数据结构都是稳定的公共接口；恰恰相反，Clang AST、LLVM 的优化流水线和机器级中间表示
都可能随版本改变。它指的是更基本的责任关系：词法阶段要保存可供语法分析的记号次序，语义
分析要把名字与声明、表达式与类型联系起来，LLVM IR 要明确控制流与数据流，目标文件要以符号
和重定位描述尚未完成的地址决定，而最终的可执行文件必须在约定输入上实现源程序所要求的可
观察行为。

这种看法也避免了另一种常见误解：教科书中的方框未必对应操作系统中的进程。Clang 的
\emph{驱动程序}（compiler driver）先构造\emph{动作图}（action graph），再把动作绑定为实际
工具与\emph{作业}（job）；预处理、编译后端和集成汇编器可以在同一个
\texttt{clang -cc1} 作业内完成。因而“预处理器、编译器、汇编器、链接器”仍是有用的职责划分，
却不能被误读为“四个必然独立运行的程序”\cite{front-clang-driver,front-clang-command}。本文同时
保留概念分层和实现现实：前者帮助推理，后者防止把一幅示意图冒充运行轨迹。

\subsection{研究对象、问题与边界}

研究对象是 Clang/LLVM 编译系统及其面向 RV64 Linux 的代码生成路径，语言侧以课程规定的
SysY 2022 子集为规范语义载体。全文只追踪一个案例：两个长度均为 8 的整型数组；程序通过
\texttt{getint} 读入前缀长度 $n$，将其限制到 $[0,8]$，逐项相乘后把每个乘积截断到
$[-8,12]$，求和并以 \texttt{putint}、\texttt{putch} 输出。这个例子小到可以逐节点、逐指令
阅读，又同时含有数组、下标、赋值、条件、循环、用户函数与外部运行时调用。它不是为了展示
“语法功能越多越好”，而是为了让同一个事实在多层表示中反复出现。

围绕这一个计算，本文回答三个问题：
\begin{enumerate}
  \item 源字符经过预处理、词法、语法和语义分析后，哪些原本隐含的事实成为带类型的结构？
  \item 结构化控制和可变状态如何依次变成 LLVM IR 的控制/数据流、RV64 指令、ELF 符号与
        重定位，最后成为进程的可观察行为？
  \item “同一程序”的判断可以得到多强的证据？编译成功、结构检查、已知答案测试与形式化证明
        分别能说明什么，又不能说明什么？
\end{enumerate}

本文的范围有意收紧。我们不把若干互不相关的小程序拼成语言特性清单，不把优化后文件变小当作
性能结论，也不声称库存 Clang 是一个符合 SysY 全部规范的前端。纯 SysY 文件是规范性
（normative）源程序；一个只增加 C 声明和预处理指令的 C 观察载体，使 Clang 能展示预处理、
记号与 AST。手写 LLVM IR 和手写 RV64 汇编是对同一语义函数的独立实现，而 Clang 生成的
IR/汇编只作为观察编译器决策的证据。这些身份若混在一起，后续所谓“等价”便没有清楚的比较
对象。

\subsection{方法：沿着一个事实走，而不是沿着工具名走}

本文在每个边界重复询问五件事：输入表示是什么；本阶段确立什么事实；输出表示是什么；案例中
哪一个具体对象承载该事实；还有什么留待下一阶段决定。表~\ref{tab:evidence-kinds} 给出全文使用
的证据标记。其目的不是给叙述增加形式感，而是避免把文档推断写成实验结果，或者把六个测试
样例写成普遍正确性证明。

\begin{table}[t]
\caption{本文的证据等级与允许的结论。}
\label{tab:evidence-kinds}
\centering
\small
\begin{tabular}{p{1.8cm}p{4.2cm}p{5.5cm}}
\hline
标记 & 来源 & 可以支持的结论 \\
\hline
观察 & 固定版本、命令与输入产生的记号、AST、IR、目标文件或运行结果 & 该环境下确实出现了所述结构或行为；不能推出所有版本都相同 \\
文档 & 语言规范、ABI、LLVM/Clang 官方手册 & 接口语义或设计责任；滚动文档仍需与安装版本对应 \\
推断 & 由观察和文档共同给出的因果解释 & 在明确假设下最合理的解释；不是工具实现对未来版本的保证 \\
验证 & 汇编/IR 验证器、结构检查和同一测试预言机 & 排除被检查到的错误；有限测试不等于全域等价证明 \\
\hline
\end{tabular}
\end{table}

这套方法刻意接近《深入理解计算机系统》的写法：不把“编译”当作一个名词，而是指着一个
对象问它是什么。第一个数组元素 \texttt{-4} 在词法层是一元负号与整数常量，在 AST 中是带
类型的一元表达式，在 LLVM 中成为 \texttt{i32 -4} 的初始化数据，在目标文件中成为四个按
目标字节序排列的字节；\texttt{getint} 从已解析的函数调用变成外部声明、未定义符号和调用重
定位，直至链接器从 SysY 运行时库中选择定义。单个对象跨层的连续身份，比罗列每个工具的全部
选项更能暴露系统设计。

\subsection{贡献与文章结构}

本文是一篇实例锚定的技术综述，不提出新的编译算法。它的贡献在于：（1）给出一个无有符号
溢出和越界歧义的 SysY 整数案例及数学契约；（2）把预处理、前端、LLVM 中端、RISC-V 后端、
汇编、ELF 链接和运行时置于同一条可检查的证据链；（3）严格区分驱动动作、实际作业与编译器
内部阶段；（4）以行为预言机和结构证据校准“等价”的强度；（5）在展望中把单层 LLVM IR 与
MLIR/AscendNPU IR 的渐进式降低联系起来，而不把未运行的硬件路径写成实验结果。

第~\ref{sec:case-study} 节固定程序的语义和旅程地图；第~\ref{sec:frontend} 节从字符走到带类型
AST；后续章节依次追踪 LLVM IR、优化与 RV64 代码生成、汇编和链接、端到端验证，并讨论 MLIR
及 AscendNPU IR。全文最终回到同一原则：抽象层不是为了隐藏事实，而是为了让恰当的事实在恰当
的阶段成为可操作对象。

\section{一个程序，一条语义线索}
\label{sec:case}

\subsection{语义契约与选择理由}

清单~\ref{lst:sysy} 是贯穿全文的 SysY 源。它读取整数 $n$，令
\[
b=\min(10,\max(0,n)),\qquad f(b)=\prod_{k=2}^{b}k,
\]
然后输出十进制的 $f(b)$、一个换行，并以状态 0 返回。故负输入输出 1，而不是 0；这是“先夹到 0，再计算 $0!=1$”的直接结果。上界 10 使结果不超过 3,628,800，落在 32 位有符号整数内，避免本例的乘法进入溢出语义争议。

\begin{lstlisting}[language=C,caption={有界迭代阶乘的 SysY 表示（注释从略）。},label={lst:sysy}]
int clamp_input(int value) {
  if (value < 0) return 0;
  if (value > 10) return 10;
  return value;
}
int factorial(int n) {
  int result = 1; int factor = 2;
  while (factor <= n) {
    result = result * factor;
    factor = factor + 1;
  }
  return result;
}
int main() {
  int input = getint();
  int bounded = clamp_input(input);
  int result = factorial(bounded);
  putint(result); putch(10); return 0;
}
\end{lstlisting}

这个程序的复杂度是有意克制的。两个有序条件分支暴露控制流；循环暴露回边和循环携带状态；三个自定义函数暴露调用边界；输入输出暴露未解析外部符号。数组、递归、浮点和堆分配会增加新的主题，却不会让本文要求的任何边界更清楚。

\subsection{两种源观察形态，而非两个算法}

预处理属于 C 翻译流程而不是 SysY 的语言契约。为观察它，配套 C 形态把 10 写作对象式宏 \code{FACTORIAL\_LIMIT}，并显式声明三个外部函数；其控制流和算法与清单~\ref{lst:sysy} 相同。此 C 文件只用于观察预处理与由 Clang 生成的产物，不另立语义。

目标契约固定为 \code{riscv64-unknown-linux-gnu}，RV64GC 指令集合与 LP64D ABI。ISA（Instruction Set Architecture）说明机器操作的意义；处理器专用 ABI（processor-specific ABI, psABI）另行规定参数寄存器、栈、ELF 重定位等接口\cite{riscv-isa-2026,riscv-psabi-1.0}。二者不可互相替代。

\begin{table}[tb]
\caption{同一事实在各观察点的预期痕迹。}
\label{tab:facts}
\centering\small
\begin{tabularx}{\linewidth}{@{}lXXX@{}}
\toprule
源事实 & 前端/IR & RV64/对象 & 最终可见结果\\
\midrule
上下界判断 & 比较、分支或选择 & 条件分支 & 输入被夹到 $[0,10]$\\
循环累乘 & CFG 回边、SSA 递推 & \code{mulw}/\code{addiw} 与跳转 & 阶乘值\\
函数调用 & 带类型 \code{call} & \code{a0}、\code{ra}、调用重定位 & 跨文件协作\\
换行 10 & \code{i32 10} 实参 & 在 \code{a0} 中物化 & 输出以换行结束\\
\bottomrule
\end{tabularx}
\end{table}

\section{从源文本到带类型程序}
\label{sec:frontend}

\subsection{预处理：改写解析器将看到的记号}

\paragraph{输入与工作。} C 观察文件包含字符、注释、\code{\#define} 指令和声明。翻译阶段把字符组织为预处理记号，删除注释并展开宏；C 标准工作草案给出了翻译阶段及记号类别的语言层契约\cite{wg14-n3220}。

\paragraph{输出与实例痕迹。} 对该文件执行仅预处理后，\code{value > FACTORIAL\_LIMIT} 和返回语句中的宏均成为字面量 10。行标记记录输出片段与原文件位置的联系；\code{getint} 等声明仍只是普通 C 声明。预处理器不会验证它们是否有实现，也不会计算阶乘。

\paragraph{边界。} 宏名作为可执行实体已经消失，但宏展开得到的记号仍须经语法和类型检查。本例没有头文件包含或条件编译，因而它们只是预处理器的一般能力，不是本产物的观测结果。SysY 源直接进入其前端，不能把 C 观察形态的宏归因给 SysY。

\subsection{词法与语法：线性记号变成源形状的树}

\paragraph{输入与工作。} 词法分析（lexical analysis）把 \code{while (factor <= n)} 分为关键字、标识符、括号和关系运算符等记号；解析（parsing）依照文法把线性序列组织成嵌套结构。Clang AST 刻意贴近源语言形状，便于诊断和源到源工具处理\cite{clang-ast-2026}。

\paragraph{输出与实例痕迹。} \code{clamp\_input} 的首个判断形成“函数声明--复合语句--IfStmt--二元比较--ReturnStmt”的层级；阶乘主体则有 WhileStmt，其条件是 \code{<=}，循环体含两个赋值。\code{result * factor} 的乘法表达式嵌在赋值右侧；\code{factorial(bounded)} 是 CallExpr，实参节点引用局部声明。

\paragraph{边界。} 树保留源码范围、运算符和语句嵌套，却不规定 RV64 指令或物理寄存器。空白、注释和宏拼写通常不再是语义树节点；“循环”此时仍是语言构造，而非基本块与回边。

\subsection{语义分析：树获得绑定、类型和值类别}

语义分析（semantic analysis）把每个标识符引用绑定到声明，检查调用形参与实参、运算类型和返回类型，并区分左值（lvalue）与转成数值后的右值（rvalue）。在本例中，\code{factorial} 的形参与返回值均为 \code{int}；\code{result} 在赋值左侧指示存储位置，在乘法中则要读取其值。\code{getint()} 的结果可初始化 \code{int input}；\code{putint(result)} 与声明的一个整数参数一致。

若把调用改为 \code{factorial(bounded, 1)}，语法仍可组成调用树，但语义分析会因实参数目不匹配而拒绝它。这一反事实只界定检查职责，不是第二个实验程序。另一方面，“$0\le b\le10$”主要由程序控制流建立；普通类型系统只知道 \code{b} 是整数，并不会自动把区间不变量编码进类型。

\begin{table}[tb]
\caption{前端三层契约。}
\centering\small
\begin{tabularx}{\linewidth}{@{}lXXX@{}}
\toprule
阶段 & 新增/显式化 & 丢弃或弱化 & 尚未决定\\
\midrule
预处理 & 宏展开后的记号、行映射 & 注释、宏作为名字 & 语法与类型是否合法\\
解析 & 语句/表达式层级 & 多数排版差异 & 名字是否绑定、目标布局\\
语义分析 & 声明绑定、类型、隐式转换 & 部分表面语法差异 & 指令、寄存器、地址\\
\bottomrule
\end{tabularx}
\end{table}

至此输出是一棵“有语言意义的树”，而不是可执行代码。下一阶段需要把嵌套语句翻译成控制流图（Control-Flow Graph, CFG）与显式操作。

\section{LLVM IR 与优化}
\label{sec:ir}

\subsection{控制流与内存效应的显式化}

LLVM IR 采用带类型的 SSA 表示，函数由基本块组成，每个基本块以控制流终结指令结束\cite{llvm-langref-2026}。在 \code{-O0} 产物中，可寻址局部变量通常对应 \code{alloca} 栈对象，初始化和赋值形成 \code{store}，表达式取值形成 \code{load}。源级 \code{while} 被降低为条件、循环体和退出基本块，循环体至条件块的分支构成回边。此处的 \code{alloca} 表示抽象栈对象，具体栈偏移仍由目标代码生成决定。

参考 LLVM IR 直接以 SSA 形式表示阶乘递推。清单~\ref{lst:phi} 中两个 $\phi$ 操作依据实际执行的前驱边选择输入定义：首次迭代从 \code{entry} 取得 2 和 1，后续迭代从 \code{body} 取得更新值。支配关系和支配边界为一般 SSA 构造中的合流位置提供理论基础\cite{cytron-et-al-ssa-1991}。案例的循环回边要求 \code{factor} 与 \code{result} 在循环头合并初始值和迭代值。

\begin{center}
\begin{minipage}{0.92\linewidth}
\begin{lstlisting}[caption={阶乘循环在参考 LLVM IR 中的 SSA 表示。},label={lst:phi}]
loop:
  %factor = phi i32 [ 2, %entry ], [ %next_factor, %body ]
  %result = phi i32 [ 1, %entry ], [ %next_result, %body ]
  %continue = icmp sle i32 %factor, %n
  br i1 %continue, label %body, label %done
body:
  %next_result = mul i32 %result, %factor
  %next_factor = add i32 %factor, 1
  br label %loop
\end{lstlisting}
\end{minipage}
\end{center}

函数调用在 IR 中同时编码返回类型、目标符号和参数类型。例如 \code{call i32 @factorial(i32 \%bounded)} 返回 32 位整数，\code{putint} 与 \code{putch} 则声明为 \code{void}。夹取条件分别采用 \code{icmp slt} 和 \code{icmp sgt}，循环条件采用 \code{icmp sle}，明确表达有符号整数比较。参考 IR 的 \code{mul i32} 未附加 \code{nsw} 标志，因而没有引入额外的无溢出假设；已声明输入路径上的乘积由上界保证有效。

\subsection{SSA 构造与标量优化}

LLVM 优化流水线由分析和变换通道组合，其默认构成随版本、目标和优化级别演化\cite{llvm-passes-2026,llvm-newpm-2026}。案例中的局部标量地址不逃逸，使标量提升能够把栈对象的读写改写为 SSA 定义和使用；常量传播、控制流简化及内联进一步改变基本块和调用边界。由此产生的文本结构差异以语义保持的程序变换为前提。

表~\ref{tab:structure} 合并报告 IR 与汇编层的结构测量。RV64 目标下，IR 指令数由 56 降至 30，\code{alloca/load/store} 从 9/13/13 降至 0/0/0，反映出局部内存状态的提升与整体结构收缩。汇编近似指令行由 77 降至 45，而助记符种类由 13 增至 15，表明文本规模和操作种类并不单调对应。这些静态指标用于刻画表示结构，性能解释不在其测量范围内。

\begin{table}[tb]
\caption{Clang 18.1.3 在 RV64 目标上的表示结构测量。}
\label{tab:structure}
\centering\small
\begin{tabular}{@{}lrrrrrr@{}}
\toprule
级别 & IR 指令 & \code{alloca} & \code{load} & \code{store} & 汇编行 & 助记符\\
\midrule
\code{-O0} & 56 & 9 & 13 & 13 & 77 & 13\\
\code{-O2} & 30 & 0 & 0 & 0 & 45 & 15\\
\bottomrule
\end{tabular}
\end{table}

优化后的 IR 仍使用不限数量的 SSA 名称和目标相对独立的整数操作。目标后端接收其数据依赖、控制流和内存语义，并将这些信息映射到有限物理寄存器、具体指令和调用约定。

\section{从 LLVM IR 到 RV64 机器接口}
\label{sec:rv64}

\subsection{合法化、指令选择与寄存器分配}

目标无关代码生成器把 IR 经过调用约定降低、指令合法化（legalization）、指令选择（instruction selection）、寄存器分配和指令调度等步骤，最终打印汇编或发射对象字节\cite{llvm-codegen-2026,llvm-globalisel-2026}。这些步骤是概念分工，不保证每个目标或优化级别都采用相同内部算法。

IR 的 \code{i32 mul} 在 RV64 上要维持 32 位结果语义。作者汇编用 \code{mulw t0,t0,t1} 计算低 32 位并符号扩展到 64 位寄存器；循环变量更新用 \code{addiw}。\code{icmp sle} 与分支则表现为 \code{bgt t1,a0,.Lfactorial\_done}：后端交换比较方向，把“若继续则进循环体”改成“若已结束则跳出”。这是同一控制谓词的目标化布局，不是语义改变。

\begin{lstlisting}[caption={作者 RV64 汇编中的阶乘函数。},label={lst:rvfact}]
factorial:
    li      t0, 1
    li      t1, 2
.Lfactorial_test:
    bgt     t1, a0, .Lfactorial_done
    mulw    t0, t0, t1
    addiw   t1, t1, 1
    j       .Lfactorial_test
.Lfactorial_done:
    mv      a0, t0
    ret
\end{lstlisting}

\code{li}、\code{mv}、\code{j} 和 \code{ret} 是汇编伪指令（pseudo-instruction）：它们改善可读性，由汇编器选成一条或多条真实编码。因此，汇编文本行数不能自动等于对象中的机器指令数。RISC-V ISA 规定真实指令效果，汇编手册和工具约定处理伪指令；本文在编码层以反汇编为准，而非把助记符外观当作硬件事实。

\subsection{RV64 psABI：独立编译仍能相互调用}

LP64D 调用约定规定整数参数和返回值优先使用 \code{a0--a7}；\code{ra} 保存返回地址；\code{sp} 是栈指针并须在过程调用边界满足对齐；\code{t0--t6} 为调用者保存临时寄存器；\code{s0--s11} 为被调用者保存\cite{riscv-psabi-1.0}。因此 \code{factorial} 从 \code{a0} 读入 $n$，在 \code{t0}/\code{t1} 中维护循环状态，再把结果移回 \code{a0}。

叶函数 \code{factorial} 不调用别的函数，可直接 \code{ret}；\code{main} 连续调用五个函数，必须保护自己的返回地址。作者汇编把栈指针减 16，在 \code{8(sp)} 保存 \code{ra}，返回前逆向恢复。这 16 字节而非任意大小来自调用边界的栈对齐约束。它也说明“有局部变量”不必然意味着“有栈槽”：这里寄存器足以承载局部状态，栈帧的主要任务是保存返回地址。

\subsection{结构观测与不能推出的结论}

同一 Clang 18.1.3 环境中，生成汇编的近似指令行由 \code{-O0} 的 77 降到 \code{-O2} 的 45，出现的助记符种类由 13 增至 15。前一变化可能来自栈访问和冗余搬运减少，后一变化说明优化会选用更多样的操作；“种类更多”不等于程序更复杂或更慢。近似计数排除标签和大多数指示符，却仍包含伪指令，因此只适于描述文本结构。

\begin{table}[tb]
\caption{RV64 汇编的结构观测，不是性能测量。}
\centering
\begin{tabular}{@{}lrrl@{}}
\toprule
级别 & 近似指令行 & 助记符种类 & 可支持的判断\\
\midrule
\code{-O0} & 77 & 13 & 保留较多显式搬运/帧操作\\
\code{-O2} & 45 & 15 & 文本更短但指令选择更多样\\
\bottomrule
\end{tabular}
\end{table}

到汇编层，类型信息已大幅弱化：\code{a0} 本身不携带“这是 SysY int”的标签；\code{mulw} 的宽度及调用点约定共同保存所需行为。符号名仍没有最终地址，而 \code{call getint} 甚至尚未决定具体位移。汇编器的任务正是把这些文本承诺编码成节、符号和重定位。

\section{ELF 目标文件、链接与装载}
\label{sec:elf}

\subsection{目标编码与重定位}

参考 RV64 汇编经汇编后形成 EM\_RISCV 类型的 ELF 可重定位对象。对象文件按节组织指令和元数据，符号表记录 \code{clamp\_input}、\code{factorial} 与 \code{main} 等已定义函数，以及 \code{getint}、\code{putint} 和 \code{putch} 三个未定义外部符号。ELF 通用规范定义节、符号和重定位的共同结构，RISC-V psABI 补充处理器相关的重定位类型\cite{elf-gabi-2025,riscv-psabi-1.0}。

对象偏移 \code{0x2a} 处对 \code{getint} 的调用具有 \code{R\_RISCV\_CALL\_PLT} 与配套 \code{R\_RISCV\_RELAX} 记录。前者关联调用位置和目标符号，后者允许链接器在最终布局已知后执行松弛（linker relaxation）。同样的结构也出现在其他外部调用处。可重定位对象已经包含目标指令字节，但跨对象调用的最终位移仍由链接阶段填写。

\subsection{符号解析、节布局与静态链接}

链接器解析输入对象与归档中的全局符号，选择所需成员，合并输入节并分配地址，随后应用重定位，生成带入口地址和程序头表的 ELF 可执行文件\cite{lld-design-2026,elf-gabi-2025}。SysY 运行时库在此提供三个输入输出符号的定义；启动代码提供 \code{\_start}，完成初始运行环境建立后调用 \code{main}。静态链接使程序执行所需代码进入同一文件。

表~\ref{tab:sizes} 给出四种入口表示产生的对象和静态可执行文件磁盘大小。C 与 SysY 对象相差 8 字节，参考 IR 和参考汇编对象依次更小；静态文件的差异仅为数百字节，因为共同的启动及输入输出代码占主要部分。文件大小同时受节、符号表、对齐和元数据影响，适用于产物结构比较。

\begin{table}[tb]
\caption{可重定位对象与静态可执行文件的磁盘大小（字节）。}
\label{tab:sizes}
\centering\small
\begin{tabular}{@{}lrr@{}}
\toprule
入口表示 & 可重定位对象 & 静态可执行文件\\
\midrule
C 观察形态 & 2,736 & 627,032\\
SysY 源 & 2,744 & 627,032\\
参考 LLVM IR & 2,312 & 626,968\\
参考 RV64 汇编 & 1,512 & 626,832\\
\bottomrule
\end{tabular}
\end{table}

链接松弛利用全局布局信息改写先前的通用地址构造或调用序列。因此，对象反汇编与最终文件反汇编可能具有不同编码，而其调用语义保持不变。

\subsection{程序头与进程映像}

Linux \code{execve} 根据 ELF 程序头表建立新进程映像。\code{PT\_LOAD} 项规定文件范围、虚拟地址、内存大小、对齐和访问权限；当 \code{p\_memsz} 大于 \code{p\_filesz} 时，额外内存区域以零初始化\cite{linux-elf-2026,linux-execve-2026}。内核完成映射和初始栈准备后，将控制转移至 ELF 入口 \code{\_start}。

\begin{table}[tb]
\caption{ELF 节视图与装载段视图的功能对应。}
\label{tab:sections-segments}
\centering\small
\begin{tabularx}{0.9\linewidth}{@{}lXl@{}}
\toprule
链接节 & 主要内容 & 典型装载段\\
\midrule
\code{.text}, \code{.rodata} & 指令与只读数据 & R--/R-X \code{PT\_LOAD}\\
\code{.data}, \code{.bss} & 初始化数据与零填充区域 & RW- \code{PT\_LOAD}\\
\code{.symtab}, \code{.strtab} & 链接和诊断元数据 & 通常不映射\\
\bottomrule
\end{tabularx}
\end{table}

节用于链接工具对文件内容分类与合并，段用于装载器建立内存映射。至此，源级函数和变量已转换为具有确定地址、权限和入口状态的机器程序。

\section{紧凑等价检查与 MLIR 展望}
\label{sec:validation}

\subsection{行为检查：支持到哪里，结论就说到哪里}

验证把“标准输出中的一个整数加换行，退出状态为 0”作为可观察契约。六个输入覆盖负数下界、零次循环、最短正输入、多次迭代、上界和超上界。表~\ref{tab:tests} 中负数先夹到 0，所以最终输出 $0!=1$；“夹到 0”不等于“直接输出 0”。C 形态与 SysY 采用同一算法，表中只把后者列作源语义代表。

\begin{table}[tb]
\caption{六个输入的期望值与三种作者表示的结果。}
\label{tab:tests}
\centering\small
\begin{tabular}{@{}rrlll@{}}
\toprule
输入 & 期望输出 & SysY 源 & 作者 IR & 作者汇编\\
\midrule
$-3$ & 1 & 通过 & 通过 & 通过\\
0 & 1 & 通过 & 通过 & 通过\\
1 & 1 & 通过 & 通过 & 通过\\
5 & 120 & 通过 & 通过 & 通过\\
10 & 3628800 & 通过 & 通过 & 通过\\
20 & 3628800 & 通过 & 通过 & 通过\\
\midrule
合计 & -- & 6/6 & 6/6 & 6/6\\
\bottomrule
\end{tabular}
\end{table}

这项检查能发现常见的边界、循环、调用和输出差错，并支持“在声明的六个输入上可观察行为一致”。它不是形式等价证明：输入空间没有穷尽，工具链和环境也是前提。更强结论需要全域符号推理或经证明的翻译验证，并明确整数、I/O 与终止语义。版本漂移也会改变打印 IR、优化流水线和指令选择；因此本文把版本、目标三元组、ISA/ABI 和计数口径列入附录，并把所有大小称为本次实验观测。

\subsection{展望：MLIR 把“何时丢弃抽象”变成设计对象}

LLVM IR 适合显式控制流、内存与机器式标量操作，但领域操作若过早拆散，后续领域变换可能无法恢复其结构。MLIR 提出可扩展方言（dialect）与多级中间表示，使领域、算法、内存和硬件抽象在同一基础设施中逐步降低\cite{lattner-et-al-mlir-2021}。本阶乘并不需要 MLIR 才能高效实现；这里只用它解释机制，不声称本地运行或加速器验证。

一个概念路径是：\code{func+arith+scf} 保留函数、整数运算和结构化循环；降低到 \code{func+arith+cf} 后，循环成为块、分支与块参数；再转换到 LLVM 方言并翻译为 LLVM IR，复用既有优化和 RV64 后端。方言转换框架用转换目标声明哪些操作合法，用重写模式替换非法操作，并可用类型转换器桥接暂时混合的类型；部分转换允许不同抽象共存\cite{mlir-dialect-conversion-2026,mlir-llvm-target-2026,mlir-toy-lowering-2026}。

\begin{center}
\small
\fbox{\code{scf} 结构化循环 $\rightarrow$ \code{cf} 块/分支 $\rightarrow$ LLVM 方言 $\rightarrow$ LLVM IR $\rightarrow$ RV64}
\end{center}

课程进阶材料点名的 AscendNPU IR VecAdd 展示了同一思想在加速器上的形态：官方 HIVM 示例把设备入口、全局内存（\code{gm}）到统一缓冲区（\code{ub}）的搬运、\code{hivm.hir.vadd} 以及写回保持为显式操作；架构说明再把 HFusion、HIVM、低层 MLIR 与 LLVM IR/算子二进制组织成逐级降低关系\cite{ascend-vecadd-2026,ascend-architecture-2026}。本文没有执行 CANN/BiSheng 专用编译器或 NPU 真机，因此这只是文档支撑的设计展望，不是本阶乘实例已经经过的阶段，也不是设备性能实验。

其价值是\emph{表示时机}：在最后一个需要高层结构的变换完成后再降低。代价同样真实：每条方言边界都需要合法性规则、重写、类型桥接、诊断和版本维护；方言也不会自动组合，保留结构更不保证程序更快。因此，MLIR 更适合确有多个领域或硬件层级、而 LLVM IR 会过早抹平关键信息的系统，不应仅因“层次更多”就用于小型标量编译器。

\section{结论}
\label{sec:conclusion}

同一个有界阶乘揭示了编译系统的连续性，也揭示了各阶段不可混同的职责。预处理器展开宏却不验证外部实现；前端把文本变成带绑定与类型的树；LLVM IR 把语句重述为基本块、效果和 SSA 数据依赖；后端在 ISA 与 ABI 约束下选择 RV64 操作和寄存器；汇编器把文本编码为字节，并以符号和重定位保存尚未知的地址；链接器兑现这些义务、安排布局和入口；加载器依据段而非节建立进程映像。

\begin{table}[tb]
\caption{端到端表示账本。}
\label{tab:ledger}
\centering\scriptsize
\begin{tabularx}{\linewidth}{@{}lXXX@{}}
\toprule
表示 & 主结构/新增事实 & 淡化或丢弃 & 留给下一层\\
\midrule
源码/预处理文本 & 名称、语句、展开记号 & 注释、宏拼写 & 文法与类型\\
带类型 AST & 树、绑定、类型、值类别 & 排版 & 控制/数据实现\\
LLVM IR & CFG、效果、SSA 定义--使用 & \code{while} 与多数源码形状 & 寄存器、指令、地址\\
RV64 汇编 & 目标操作、ABI 动作、标签 & SSA 名与多数类型 & 编码与符号地址\\
可重定位 ELF & 节、符号、重定位、字节 & 指令文本外观 & 全局解析与布局\\
可执行 ELF & 入口、程序头、已解析布局 & 多数链接义务 & 映射与初始状态\\
进程映像 & 地址、权限、机器状态 & 多数节名和符号名 & 实际执行\\
\bottomrule
\end{tabularx}
\end{table}

表~\ref{tab:ledger} 也修正“降低只会丢信息”的直觉：每层都删除对下一任务无用的区别，同时创造新的、可检查的事实。最终进程里没有宏、WhileStmt 或 \code{\%result}；但它们要求的行为存活在有类型的运算、数据依赖、目标约定、重定位后的位字段和加载映射中。六个测试支持这条具体链的行为一致，计数则只帮助看见结构变化。真正“理解编译系统”，不是背诵命令，而是能在每个边界指出：现在的表示知道什么，为什么知道，以及下一消费者仍需要解决什么。

\appendix
\renewcommand{\theHsection}{appendix.\Alph{section}}
\section{可复现实验清单}
\label{app:reproducibility}

提交前应由自动化脚本生成或核对以下材料：
\begin{enumerate}
  \item \textbf{身份与环境：}操作系统、体系结构、容器或虚拟机信息；编译器、LLVM、链接器、模拟器及 MLIR 工具的完整版本输出。
  \item \textbf{输入：}SysY、LLVM IR、RISC-V 汇编源文件，运行时库来源与校验和，全部测试夹具及其预期输出来源。
  \item \textbf{命令：}预处理、IR 生成、优化、汇编、链接、反汇编与执行的可复制命令；不得用省略号替代关键参数。
  \item \textbf{制品：}预处理文件、未优化/优化 IR、汇编、目标文件、符号表、重定位表、链接映射、控制流图和逐次运行日志。
  \item \textbf{判定：}差分脚本版本、输出规范化规则、失败样本保存策略和汇总文件的生成方式。
  \item \textbf{MLIR：}VecAdd 来源版本、环境安装步骤、逐阶段通道命令、各阶段 IR 与失败诊断；设备不可用时明确可复现边界。
\end{enumerate}

推荐从仓库根目录使用版本化脚本完成实验，并将临时文件放入项目内的 \texttt{.temp/}、缓存放入 \texttt{.cache/}。论文构建本身见 \texttt{report/README.md}。最终 PDF 应由干净检出（clean checkout）生成，以避免本地未跟踪字体、宏包或制品造成隐式依赖。

\bibliographystyle{splncs04}
\bibliography{references}
\end{document}
